\section[第一章\q 走向交易者之路]{第一章\\ 走向交易者之路}

\subsection{学生股民}

{\kaishu
也许你会赢，也许你会输，但无论如何你都会失去很多你最最宝贵的东西。

\hfill ——电视剧《坐庄》中的一句对白

做交易是找死，不做交易是等死。

\hfill ——金融帝国
}

过去两年波澜壮阔的大牛市，使得越来越多的人开始关注股票市场。其中也不乏一些学生，他们也梦想着通过资本市场来实现自身的价值，改变自己的命运。当我第一次听到“学生股民”这个词时真的倍感亲切，毕竟十多年前我就是一位名副其实的“中学生股民”。

前几天，我的一位好朋友对着他女朋友“抱怨”说，他都是跟我学坏的。他说上学的时候，我教他抽烟、喝酒、赌博。出人意料的是，他女朋友反问道：“他怎么没有教你炒期货呀？”他很经典、很搞笑地回答：“谭总（我）说了，这个可不能学。”

也许是经受的挫折太多了吧，我真的从来没有劝过身边的人去做交易。就连最近很多同龄人让我教他们炒股票时，我也都是尽量地推托，并且和他们讲明炒股的弊端。因为我深知一位交易者的成长是非常痛苦的，这需要和行业环境、自身本能、信念体系等一切固有的东西对抗。我从来不相信有人能够不经历惨重亏损，就能真正学会做交易。事实上，初学者赚来的钱，通常都是市场借给你的，等到市场向你讨账的时候会更加凶狠。似乎这个世界上什么东西都能教，但唯独两件事情只能靠自己去感受。一个是恋爱，另一个就是做交易。你无法向一个儿童解释恋爱的感觉，更无法教会他恋爱的技巧。事实上，交易也是别无二致的。

我深信我说出来的东西，初学者是不会理解的；就算是理解也不会相信；就算是相信也不会去实施；就算是实施也不会坚持；就算是坚持也绝对不可能在没有经历过惨重亏损以前，就能心甘情愿地坚持到底！事实上，交易真的并不复杂，相反却是出奇的简单。问题是这种简单有些出乎常理，并且不符合正常人的本能和信念体系。本书的全部内容可以浓缩为几句话，并且我保证这些话大家肯定都听说过。其实本书的全部价值就是让大家能够相信这些简单的道理。

我在网上看过一篇关于进入股市动机的调查。排在第一位的是获利；第二位的是证实能力；第三位的是追求自由。回想一下，我当初也确实因为这些才进入股票市场。但具有讽刺意义的是，我首先是付出了高昂的学费；然后是必须愚钝地坚持交易原则；最后的结果却是基本上靠天吃饭，而决无自由与自豪可言。十年后的今天，我不得不感叹，结局与初衷真的是事与愿违！我相信现在进入股市的这些同龄人，也许有一天你们真的能够赚到庞大的财富，但你们肯定会感慨市场和你们当初想象的并不一样。

当我回想起十多年来经历的一幕幕：那无数的不眠之夜；那种面对亏损绞尽脑汁也找不到出路的无奈；那种毁掉爸妈一辈子积蓄三分之二的心痛；那种与本能较量的痛苦挣扎；那种为了交易而离群的孤独……现在心里还是酸酸的。说实话，我真不忍心看着身边的同龄人再去经历我曾经历过的那种痛苦的折磨，要知道他们还只不过是一些孩子。我更不忍心看着这些充满童真的可爱心灵，变得千疮百孔后生出铁石一般的老茧。即使这样，谁也不能保证这种付出必然能够有好的回报。

最近的一些事情让我的观点有所转变。我发现身边的很多同龄人也都情绪低落且充满压力。房价在飞涨，我们这些二十多岁正要告别童年的孩子真的是有很大的压力。他们也都在各自不同的行业努力为钱而拼搏，也许任何行业都不可能一帆风顺吧。有一位心理学家说过，大众通常都会觉得自己付出的心血比其他人更多，毕竟只有自己更能感受到自己的付出。以前是“人无外财不富”，也许现在是“人无外财难活”。毕竟没有什么行业能够像交易这样有如此低的进入门槛，又有如此惊人的短期回报。

我算了一笔账，我老妈像我这么大的时候一个月工资不到40元，而现在一套40平方米的二手房就要将近40万。要知道这只是非常低的住房要求，好一点的房子动辄上百万，而现在的平均工资估计在2000元左右。是否等我们这些“80后”步入中年的时候，房价要涨到2000万呢？客观地说，用过去20年来推导未来20年，似乎并不存在太大的可比性。但谁也不能否认通货膨胀是经济发展的必然产物。考虑到复利效应，即使是很温和的通胀，在20年以后也是非常可怕的。在全球资本过剩的今天，虽然物价飞涨，但利率还是低得可怜。如何理财已经是我们这一代人不容忽视的问题！

\subsection{皇帝的新装}

{\kaishu
能够为你钱袋负责的人只有你自己。二选一，独立担负起责任，还是离开？

\hfill ——金融帝国
}

五年前我离开了有形的实体市场，在家中通过虚拟网络完成非现场交易。最近出于一些原因，我又回到了阔别五年的现实市场中。这五年来我养成了一种与以前截然不同的交易理念，本以为现在的这个行业也会有很大的改变，但遗憾的是，这个行业与五年前根本没有任何不同。

在我看来，这个行业中的每一个人都在说着一些或许幼稚、或许言不由衷的话，并且各自扮演着不同的角色，共同演出这场盛大的闹剧。我希望本书能够像那个童话中天真无邪的孩子一样，告诉大家皇帝根本就没有穿衣服！

在期货市场的交易大厅中，我听到最多的一句话就是：注意隔夜风险，建议短单操作。所谓期货营业厅中的研发部，就是在做着为了自身利益而向交易者灌输有利于增加成交量的思维方式。他们通常是大学金融系刚毕业的学生。虽然没有任何实盘经验，但却熟悉经济学理论。自然他们就以一种固有的方式来解释行情，并让初学者们感到他们的博学。基本来说，他们就是用最快的速度搜集所有基本面的信息，然后告诉你上涨或者下跌是受什么消息的影响。最后再说一些相反的消息，告诉你注意风险，不建议过多持有隔夜单。当然如果没有什么值得一提的消息，那就说是受技术性买（卖）盘推动（打压）……偶尔他们也会谈谈支撑位和阻力位。有意思的是，他们的说法和我五年前在股市上听到的股评惊人相似。那时大盘在一千多点，对于股评而言，就是向上看50个点是阻力，向下看50个点是支撑。无论是在牛市还是在熊市中，似乎他们认为大盘就应该静止不动。除非是大盘运行到顶部或者底部，否则他们是不会意识到趋势的存在。要知道市场总是在发威的时候比最有想象力的预测者更加具有想象力！有人说期货的获利方法不外乎是：用小的仓位来赚大的波动，或者是用大的仓位来赚小波动。这话说得一点不错，期货交易者通常会找到风险与仓位的平衡点。似乎期货公司最不乐意看到交易者用轻仓来应对“巨大”的隔夜风险，自然他们更乐意把交易者引向重仓与短线的深渊。

在一个能容纳数百人的四星级宾馆的大厅内，我听到几位资深分析师举行的一场大型报告会。我个人感觉，这种以招揽客户为目的的报告会，一定会给潜在客户一些最好的服务。事实上，这时期货公司最大的利益不是什么手续费，而是得到大众的认同。分析师讲道，趋势是你的朋友，逆势交易是最可怕的事情。所以牛市只能做多，熊市只能做空。牛市逢回调买入，只会买贵、不会买错。只有顶部能够做错一次，而之前你能做对N$-$1次……完全正确，讲得很到位，这样冠冕堂皇的话，让我都迷惑了几分钟。请注意，“只会买贵，不会买错”是什么意思？我的理解是，亏损了不要走，即使跌了也只是贵了而不是错了。天哪！这是和趋势做朋友吗？典型的逆势交易！趋势的朋友应该是：在牛市中持有多单，每逢超过一定幅度的调整时平多出来，等到整理结束后再以更高的价位买入。资深分析师？居然没有搞清楚正确率与赔率的问题吗？典型的“求势不舍价”的草包思路！

此后我又听到了很多看似冠冕堂皇，实质非常荒唐的说法。这种报告会对交易者更具有杀伤力，他们用对真理的曲解来误导交易者。最后，我收到了几份打印出来的研发报告。理由、论据都很充分，观点也很明确，保护性止损也被考虑到其中。但有一个致命的缺陷：基本上都是理论上可行，而可操作性不大。报告中所有内容都在试图证明自己比市场聪明，而对市场未来的走势指手画脚。也许分析师只能是一些能够看懂行情与市场为敌，并且最终能够战胜市场的人吧。

当读者看到这里时，也许会对这个“乌烟瘴气”的行业愤愤不满。但大家可曾想到过深层次的问题？这到底是谁的错呢？我个人觉得无论是研发部还是分析师都没有错，错都在交易者自己身上。想想看，当你去商场购物的时候，哪一个商家不是说自己的商品质量好、价位低？你见过哪个商家告诉你某某产品的质量低于行业平均水平？你见过哪个商家告诉你什么地方还有比这里售价更低的商店？商家都是逐利的，自然所有的宣传都是利己的宣传。然而，人通常就是这样，在社会上聪明，可一进入市场就糊涂起来。期货公司与商家并没什么不同，他们既不是你的朋友，也不是你的敌人。确切地说，你们应该算是博弈关系。如果没有商店你去哪里购物？如果没有期货公司你又去哪里下单？消费者绝对不会把自己的利益寄托在商家身上，但交易者却经常把自己的利益寄托在期货公司身上。

几年前我在朋友的一家手机店做过一个试验。我对消费者讲解的时候，以我全部的专业知识向他们介绍各款手机的优势和弊端，希望他们能够结合自身的需求来进行最优选择。遗憾的是，消费者并没有被我的真诚而感动，相反他们更乐意从“油嘴滑舌”的店员手中购物。要知道，没有弊端的完美手机（质量最好、功能最全、款式最新、价格最低）并不存在，但消费者想要的就是那种完美的幻觉，所以那些靠销量为生的店员就必须成为制造幻觉的高手。

再和大家讲一个故事吧。我曾经听到一位朋友向广告公司抱怨KT板的质量越来越差的时候，说出了这样几句话：“你们知道这是为什么吗？因为消费者相对于质量而言更注重价位。如果消费者更乐意买入价格高而质量好的产品时，产品的质量就会越来越好。当你和广告公司讨价还价的时候，广告公司也会和生产厂家讨价还价，最终厂家只好降低成本。”这说明，永远是需求在决定供给。

大粪和花朵谁会更有市场？这取决于要“招”的是苍蝇还是蜜蜂。我的第一份工作是在手机制造厂家做销售代表，工作内容就是开发和维护与经销商的渠道。我的工作经验告诉我，站在客户的角度上考虑问题，并且诚信地开展业务就是做渠道工作的法宝。手机行业没落以后，我不得不考虑改行的问题。既然我从小就做股票，那么为什么不试试金融行业呢？由于我没有学历，所以只好从最基础的工作干起。遗憾的是，我用做“渠道”的思路来从事“推销”性质的工作，根本就是行不通的方法。在银行大厅一位中年妇女向我询问，买什么基金比较好？我回答道：不存在最好的基金，而是应该挑选一只适合你的基金。你是准备长期持有，还是只想管理一下流动资金呢？你是希望绝对不会亏损，还是可以承担一定程度的亏损，并且追求更高的收益呢？那位中年妇女马上就急了！近乎愤怒地跟我说：“没有你这么说话的，我投资是为了赚钱。难道我买基金就是为了亏损吗？”这是发生在大牛市顶部阶段的一件小事，或许大家无法想象：别说是预测市场会下跌了，就是告诉人们买股票是可能存在风险的，都是一件很难让人接受的事情。

我突然意识到了做渠道与搞推销之间的根本区别。对于做渠道而言，经销商既然能够自己开店当老板，一定是比我更懂手机的人。甚至他接触过形形色色不同品牌的销售代表，我是否“实在”人家一眼就能够看得出来。这就是和强者做生意的好处，复杂的生意可以用最简单的方式来谈。但与弱者做生意就明显不同了，他们比你懂得要少得多，但同时他们又是你的上帝。很显然，申购货币市场基金我是不可能赚到钱的，而申购股票型的开放式基金却可以让我赚到不少的佣金。事实上，客户经理存在的意义就是在银行大厅推销能够赚取高额佣金的开放式基金。但真正了解金融市场的人，就算想买入股票基金也只会买入封闭式基金，毕竟它的手续费不足前者的十分之一。我甚至尝试过放下自己的个人得失，来推销封闭式基金。遗憾的是，几乎没有人能够了解我的良苦用心。后面大家将会看到，与此类似的策略让我在做渠道的工作中，取得了多么不可思议的成功。

这段工作经历不得不让我进行一次痛苦的抉择：坚持还是离开？很显然，我不可能一边告诉客户大熊市就在眼前，同时又能够完成公司交给我的开户任务。我记得公司经理跟我说过一句很经典的话：你不要考虑太多，这个行业存在了这么多年，用什么方式、说什么话可以开发最多的客户，公司一定比你清楚。言外之意就是，公司让我怎么跟客户忽悠，我就怎么跟客户忽悠。难道融入金融行业就是让我放弃良知与自由意志吗？让我更加心痛的是，即使我为客户着想，客户也不可能知道、更不可能念我的好。在我看来，良知与自由意志是最为珍贵的东西，因此，最终我放弃了进入金融行业的想法。当然，我并不是说我的同事都是没有良知的人。最好的推销者绝对不是没有良知的人，而是能够先把自己骗了然后再去骗别人的人。毕竟没有人喜欢永远带着面具且言不由衷，恐怕这就是公司经常开会的原因吧。我当时的经理也承认：我从事工作的障碍不是懂得太少了，而是懂得太多了。而那些大学刚毕业、涉世未深的小孩，则更容易先把自己骗了。

分析师的所作所为其实很简单，不过是为了迎合交易者的需求，确切地说，是在迎合人性。也许真正资深的分析师，也曾遇到过我上述的那种种矛盾吧。最终的结果只能是给客户他们想要的东西，而不是他们需要的东西。交易者喜欢做高难度动作，那么分析师就帮你尝试；交易者喜欢占市场的便宜，那么分析师同样帮助你；交易者希望次次正确，那么分析师只能宣传“战无不胜”。在这个行业中生存有两条路可走：要么就做研发、要么就做实盘。为什么两者就不能兼容呢？答案很简单：交易者是敢于承认自己看不懂行情的人，所以才会出现什么止损点；分析者是必须宣称自己看得懂行情的人，否则他们就失去了在台上滔滔不绝的价值。最可悲的就是想取两者之长的人，结果只能是不伦不类的荒唐。比如，找到10条基本面向好的理由，然后设置一个只有百分之几波动幅度的止损点。

在我看来，对于大多数中小投资者而言，趋势跟踪是一种最好的策略。几乎每一本关于交易的书籍都会或多或少地介绍趋势的重要性。但有一个不容忽视的问题：很少有书籍能够阐述趋势跟踪的副作用。如果你采用趋势跟踪策略，那么基本上就等于告别了高正确率。任何交易方法都是有其优势与弊端的。遗憾的是，整个行业中的全部角色都会向你宣传优势，而对弊端进行绝对的回避。如果交易者不能对一种方法的弊端有足够心理准备的话，那么就几乎不可能把这种方法坚持下来，自然也就绝对不可能体验到这种方法的价值。大量的交易者都是在相信顺势而为能够获利，但却并不清楚会经常出现“买在天花板，卖在地板价”的成本时，而贸然进行趋势跟踪。几次让你无法接受的“荒唐”亏损后，你就会离趋势跟踪思维而去。你付出了成本，但却没有足够的信念等到丰收的那一天。《亚当理论》是对“顺势而为”阐述最为深刻的书籍之一，但书中对低正确率的弊端同样丝毫没有提及。

1998年我就对止损策略的好处深信不疑，这要感谢许多交易书籍对我的启发。但遗憾的是，这些书籍只是告诉我止损的好处，而没有让我知晓止损也是可能无法避免地出现连续无效的事实。2000年的牛市运行得非常复杂，其间出现了很多次有惊无险地快速大幅回落。我在每一次大幅回落时都采用了止损策略来保护自己，但市场总是一次次地嘲笑我。身边的人都因大胆地逢低吸纳而一次次地占到便宜，这给我带来了巨大的压力。三次无效止损后，我终于放弃了止损策略，结果就是那唯一一次没有止损的交易，让我亏损了6位数的资金。这就好像一位医生给病人开具了一付药方后就离病人而去。医生并没有告诉病人此药需要服用多长时间才能见效，也没有告诉病人在见效以前会有什么样的副作用出现。结果病人在服药后，不但丝毫没有见到病情的好转，反而又出现了众多可怕的不良反应。这时病人已经无法再向医生咨询，怀疑与恐惧使得全部治疗最终前功尽弃。君子算战先算败，一味以获胜后的利益来怂恿君主开战的军师绝对不是什么好军师。

与此有异曲同工之妙的是，最近几年才出现的某知名分析软件。它那“B点买入，S点卖出”的指标确实直观有效。从本质上来说它与SAR指标并没有任何不同，也同样属于趋向指标的范畴。我相信如果交易者能够坚持使用这套系统，数年后绝对能够赚取不小的财富。但这套交易系统从宏观上来看，必然也是正确率低于50％的系统。这套软件在销售中一边夸大了系统的神奇（通过展示最优案例），一边又对低正确率的事实只字不提。那么我相信，使用这套系统的交易者，很可能在经历了几次左右挨打后，将能够大赚一笔的买入信号忽略，或者因忽略卖出信号而惨遭深度套牢。最终的结果就是，系统测试能够获利，但使用的交易者却出现亏损。由此看来，获得能够让你获利的交易技术，要比获得能够让你获利的交易信念简单得多。还有一点挺有意思，对于趋势跟踪来说，买入点应该在阳线的上方，卖出点应该在阴线的下方。但该软件买卖点的标注正好相反，从而使交易者在视觉上产生一种可以赚得更多的错觉。

其实这些只不过是这个行业中误区的一小部分，更多的误区我会在后面慢慢阐述。比如说，交易竞赛比拼收益率等等。这些都只能让交易者离最终的成熟越来越远。这里点到为止地讲出这些，只不过是告诉大家：在这个市场上唯一能够信任、能够依赖的人只有你自己。如果你真的想成为一名交易者，那么就必须担负起为自己的出路和钱袋负责的使命。请相信我，这个行业中的每一种角色都有自身的需求，而绝对不会有人真心地为你着想。

下面就是本章的最后一节了，我只阐述一些实事，但先不告诉大家这些现象说明的道理。对于初学者而言，我希望你们能够对交易之路的坎坷做好充足的心理准备。如果这些都不能吓退你，那么我由衷地祝愿你们未来能有好运，并且支持你们所作出的决定。对于老手而言，我会一点一点地带领你们探讨解决之道。

\subsection{大师的结局}

{\kaishu
为了保护自己的信念，我们只能对一小部分真相视而不见。

\hfill ——金融帝国
}

\textbf{本杰明·格雷厄姆（Benjamin Graham）}

格雷厄姆是美国华尔街的传奇人物，被称为“现代证券之父”。他的著作《证券分析》与《聪明的投资者》在西方享有极高的声誉，被称作“投资者的圣经”。格雷厄姆开创了价值投资的先河，其思想启发和影响了一大批后来的顶尖证券投资专家，其中巴菲特就是他的得意门生。

在我看来，格雷厄姆只不过是开创了价值投资的先河，而对于技术分析与交易策略等并没有什么太大的贡献。所以将其称为“现代证券之父”似乎有些言过其实，但将其称为“基本分析之父”绝对一点也不过分。

我最喜欢他的一句名言是：聪明的人用不需要聪明脑袋的方法赚钱，不聪明的人用需要聪明脑袋的方法赚钱。格雷厄姆还有一段话，能够清晰地反映出他的理念框架：根据一个古老的传说，一个聪明人将世间的事情压缩成一句话，“一切都如过眼云烟”。面临同样的挑战，我大胆地将成功投资的秘密归结为四个字，那就是安全边际。由此看来，格雷厄姆理论似乎是沿袭着“价格围绕价值上下波动”的理论而发展起来的分析方法。至少在他年轻的时候，他认为自己能够预知市场的未来，能够指出市场应该如何运动。所谓安全边际不过就是占市场的便宜。

这位华尔街的教父于1976年去世前不久，在接受美国《财务分析师》杂志（Financial Analysts Journal）的访谈中，却宣布他不再信奉基本分析流派，而最终相信“效率市场”理论。他认为，靠证券分析方法中刻意创立的分析技术，已不再能发现超值获利的投资机会。在他的《有价证券分析》出版的年代，确实存在这样的机会，但是现在情况已经大大不同了，当整个投资行业都在用同样的方式来发掘超值股票时，分析的成本就极大地提高了。因此，格雷厄姆最后宣称，他最终站在了“效率市场”理论一边。

\textbf{威尔德（J. W. Wilder）}

威尔德早期是技术分析指标的推崇者，著名的相对强弱指标（RSI）就是他发明的，他还发现一些其他技术指标或原理，如动力指标（MOM），市价波幅和抛物线指标等，当时受到不少投资者的欢迎。强弱指标（RSI）沿用至今仍是被广大投资者重视的分析技术指标之一。

我手上没有更多的历史资料，所以无法得知威尔德更多的个人经历。RSI指标应该是最为古老的技术指标之一。上个世纪90年代，我国的钱龙软件沿用台湾的技术指标中就有很多是威尔德的发明。虽然现在由于计算机技术的普及，技术指标已经越来越多地出现在大家面前，但在我国股票市场建立初期，技术指标在钱龙软件中的数量并非很多。在我看来，威尔德即使不能算上“技术指标”之父，也肯定是在技术指标形成初期做出重大贡献的人。

奇怪的是，威尔德最后发表的一些文章中，推翻了这些技术分析工具的好处，提出了当时崭新的理论——亚当理论来取代这些辅助技术分析指标。在《亚当理论》的著作中，我们可以反复地发现他对技术指标的评价为“主观和充满缺陷”，而宣扬顺势而为、保持弹性、敬畏市场的先进理念。后来我才知道，他嘲笑的那种自作聪明、主观地指出市场应该如何去做的想法，其实是在嘲笑自己前半生全部努力的结晶。

\textbf{杰西·利物默（Jesse Livermore）}

《股票作手回忆录》的主人公，他对市场本质认识的深刻程度让我感到惊讶。书中无数经典的语句在近百年后的今天还在被后人反复引用。不夸张地说，我就是希望靠交易理念制胜的人。本人交易理念的形成很大程度上受《通向金融王国的自由之路》的影响，该书很细致地讲述了系统交易的理念。如果寻找源头的话，可以找到更简洁并且接近本质的读物《克罗谈投资策略》。如果继续沿着理念之路寻找下去，我们可以找到更古老的《股票作手回忆录》。至此似乎已经无法找到更早关于市场本质、交易理念的著作了。所以说，将利物默称为“交易理念之父”一点都不过分。要知道利物默已去世超过半个世纪，以他的生平写的书《一位炒家的回忆》出版于1923年，几十年来再版无数次，今天还是每位炒手的必读书。

我无法否认杰西·利物默是一位天才，同时也是包括我在内众多交易者的信仰所在。我曾经在和一位网友探讨是否应该绝对地放弃思考，而机械地随势而动的时候，那位网友建议我看一下《股票作手回忆录》。有意思的是，我打算绝对放弃思考从而绝对地敬畏市场，也同样是受到了杰西·利物默的启发。在我看来，其他大师或多或少都有一些理论研究的味道，而杰西·利物默理念中展现出的智慧火花完全是来自于实战！那位网友并不知道杰希·利物默的结局，他的结局确实是所有投机者最最不愿意相信的真相。

杰西·利物默曾经用微薄的资金创造出天文数字的利润。他曾经破产过两次，但不久后他都能创造出更大的辉煌。但当他第三次破产时，他选择了自杀。他离开这个世界的时候，只留下了一千多美元的遗产。这就是杰希·利物默的结局。

很多交易者都在不停地探讨利物默陨落最本质的原因。《十年一梦》的作者就是在深受利物默启发后，无奈地接受其死于自杀的事实。利物默曾经深刻地指出，期货交易的三大陷阱为：逆势、重仓、短线。理论上讲，如果他没有违反其定下的三条戒律的话，他有可能失败、也有可能大幅亏损，但绝对不会走向灭亡。事实上利物默曾经经历过三次破产。前两次破产后都能继续用微弱的资金东山再起，并且再创快速致富的神话。

也许他还能够第三次东山再起，但他却用他的死来向一个世纪以后的我们传递一个重要的信息：你也许能够成为一时的国王，但你永远无法击败市场。他用自己的死向市场表示屈服、向市场表示敬畏。前两次破产是财富的崩塌，而最后一次应该是信念的崩塌。我相信他是死于绝望而不是贫穷。也许他一生的梦想就是击败市场，而最终他用最彻底的方式向市场投降。也许在利物默看来，他有信心第三次东山再起，但他仍旧逃脱不了第四次破产的命运。如果他人生的使命就是战胜市场，那么在他看透了市场的本质后选择了离开这个世界。

让我在此向利物默表示最崇高的敬意。他的一生是成功的，至少他最终明白了市场是不可战胜的事实。遗憾的是，近一个世纪后的我们还在试图否定天才用死来传递的信息。

\textbf{艾略特（Ralph Nelson Elliot）}

艾略特与查尔斯·亨利·道是同一类人，他们应该算得上“市场研发部门之父”。预测与交易也许本身就应该使用两种截然不同的思维模式。利物默是一个典型的交易者，而艾略特是一个典型的研发者。艾略特发明了著名的波浪理论，但我似乎没有看到任何关于艾略特本人利用波浪理论获利丰厚的报道。的确，艾略特绝对是一位应该令人尊敬的研究者。

艾略特本人是一位专业会计师，他卧病在床后致力于研究股市，研究反映市场运动的道琼斯指数。经过长期的研究，发现股市也有相同的重复现象，并提出著名的波动原理。而查尔斯·亨利·道则是道琼斯公司的创始人，并且担任《华尔街日报》的编辑。同样发表道氏理论的查尔斯·亨利·道，我也只是看到了他完美预测的报道，而同样没能找到其获利丰厚的记载。

我至今仍然认为波浪理论是一个非常完美的理论。甚至我相信，如果有一天我想将自己的交易系统升级到一个更高的层次的话，那么波浪理论很可能是一把打开大门的钥匙。但请大家思考一个问题：艾略特为何没有通过他的理论创造出财富神话呢？本书其后会进行详细的探讨，这里略微提示一下：也许波浪理论是一个可以预知未来的方程式，但可惜它无法得出一个唯一确定的解。

\textbf{江恩（William D. Gann）}

江恩也许是最有争议的人物。有人说他在市场上找到了看穿未来的方法，从而赚到了天文数字的财富。同样也有人说江恩只不过是一个骗子，他靠包装自己来获取高额的培训费。据江恩的后代所说，江恩根本没有给他们留下多少遗产。也许江恩注定永远是一个神秘的人物，毕竟关于他的历史记载充满矛盾。

在我看来，江恩理论确实是非常有价值的。我无法否认价格运行速度确实具有倍数关系的特征；同样也无法否认江恩百分比回调能够经常创造出的和谐；更无法否认时空转换的完美。但说实话，什么江恩轮中轮、时间周期等等，我真的没有弄懂。七八年前江恩真的是我最崇拜的人，我曾长时间研究他的理论。那时的动机和大家并没有什么不同，也是希望找到看穿市场未来的方法。事实上，在对历史行情的解释上，我所知的江恩理论几乎是经常创造出完美的吻合。如果有人说江恩曾经在市场上赚到了天文数字的财富，那么我绝对不会有丝毫的怀疑。

在我看来，江恩确实也找到了预测未来的公式。但遗憾的是，这个能够预测未来的公式，同样根本无法给出一个确定的解。就像爱因斯坦所说的：能够预测未来的公式中必然存在不确定性。这说明什么呢？这说明江恩理论用于解释行情能够带来完美，但应用于实盘交易就存在很大的障碍。我曾经在过去的数年间用江恩理论和波浪理论找到了那种传说中的完美与和谐，但我始终无法将这种完美用于实战。

在我看来，江恩同利物默有一点本质上的不同。无论《股票作手回忆录》有多少细节的描述，最后都能归纳为一句流芳百世的名言。而江恩似乎真的有刻意制造神秘感，而故弄玄虚的味道。如果你看《亚当理论》，你能够感受到一种作者转递知识的欲望。本书可以归纳为“顺势而为”四个字，但作者为了让读者相信这四个字，能够用大量篇幅说得透彻。而江恩似乎给人一种欲言又止的味道，并且勾起读者探索未知的欲望。这一点在江恩的著作《时空隧道》中可以反复地感受到。所以说，如果有人认为江恩是一个善于制造神秘与幻觉的魔术师，从而大赚观众钞票的人，那么我也绝对不会反对。

在我看来，江恩同利物默还有一个非常奇特的共同点。那就是意识中对市场无比的敬畏，而潜意识中又有一种无法克制的征服市场的欲望。江恩《华尔街四十五年》中提到最多的就是“止损点”这三个字。同时似乎在向我们传递一个有意思的提示：对于我这样的预测高手来说，规则的重要性都要大大高于预测。从“计划你的交易，交易你的计划”中可以清晰地感受到一种对市场策略化的敬畏。说到江恩对市场征服的欲望，我想大家都会有很大的同感。其实能够把这种欲望描绘得淋漓尽致的是下面1909年《Ticker》杂志上的一段记载：

{\kaishu
1909年夏季，江恩预测九月份小麦将会是1.20美元。但9月30日，芝加哥时间12点，小麦的价格还低于1.08美元。在预测眼看就要落空时，江恩自信地说：“只有在收市时仍未达到1.20美元时才能证明我的预测方法有误，我并不关心现在的价格，我认为价格一定会达到1.20美元。”最终1909年9月30日的小麦收盘价令人震惊地升至1.20美元……
}

我不否认这是曾经最令我激动的梦想，但现在想想看，事情的真相也许是这样的：我只学到了江恩理论的皮毛，也能身临其境地感受到那种和谐与完美。我相信这段记录很可能是真实可信的，江恩或许有能力作出这样完美的预测。但我更加深信这只不过是江恩理论给出的几个“解”中的一个。否则，这就与所有的现代哲学与科学的论点相矛盾。如果这种“解”只有一个，并且具有再现性的话，那么江恩对人类社会的推动将超过爱因斯坦。也许这就是历史上所有的奇人都无法经常再现奇迹，从而逃脱不了“善卜者不问”托辞的原因。

从江恩身上我看到了一位理论大师创造出的和谐与完美；也看到了只有交易大师才能得出的市场感悟；同时又看到了一位销售大师创造出的魔术师般幻觉。据说，江恩所有著作的版权都不属于他的儿子，这只能有两个原因：其一，江恩临终前对自己所有的著作都不真正赞赏，他不想传给子孙，怕误导他们，在《华尔街四十五年》中已看不到他年轻时发明的神奇理论的半点影子，而只是用极大的篇幅来记录行情数据，这便是例证；其二，晚年的江恩已经穷困潦倒，他不得不出卖自己所有著作的版权，这一点从记者对他儿子的采访中可以看出，他儿子说江恩的遗产只有10万美元。也许这些推论并不完全可信，但也绝对有其可信的地方。至少两点是有证可查的：第一，江恩给儿子留下的遗产只有10万美元；第二，江恩所有著作的版权都不属于他的儿子。

\subsection{总结}

我们可以从这些现象中发现什么问题呢？上面这些大师包含了基本分析、技术指标、神奇理论、交易理念的创始人，应该说他们都处在各自专业领域金字塔的顶峰，甚至有的还是开山祖师。而他们的结局又怎么样呢？有的在晚年推翻了自己前半生的努力和信念；有的最终被市场彻底打败；有的只能作理论研究而无法进行实战。大师就是大师，他们是值得敬佩的！要知道否定自己是多么困难的事情。很多人就是因为一个错误的起点，而不得不走向可怕的深渊。在科学界、哲学界，甚至是政治界，又有多少伟人能够最终做出自己否定自己的事情呢？也许交易之路就是一条否定之路，也许突然我们有一天会发现，之前我们的全部都是大错特错的。

具有讽刺意义的是，被发明者否定、嘲笑的RSI指标，在这个世纪还被众多交易者视为法宝，并且至今仍旧保留在几乎所有的行情软件中。当然我的交易理念如果向前追溯，支持理念基石的创始者居然死于破产后的自杀。似乎交易本来就是一个悲观的行业，初学者们，你们真的想好了吗？本来我还想讲述一下，住友财团的滨中泰男、巴林银行的尼克里森、中航油的陈久霖，是如何走向毁灭的。其实这方面的介绍网上有很多，大家可以自己看看。

当然，我相信在这个具有上百年历史的市场上，可能有更多的人通过市场创造出足以实现财务自由的财富，并且过着幸福快乐的生活。只不过没有人知道他们的名字，也许他们从来都没有辉煌过。这是我第三次、也是最后一次的警告——你想好了吗？如果你还想去走这条充满坎坷的交易之路，那么本书试图剖析大师结局的深层原因，但愿能够给你带来哪怕一丝的启发。

\clearpage
